\section{Discussion and Limitations}
The construction supplies a distribution-level justification for a geometric midpoint loss and a direct generator algorithm. Its guarantee requires finite-energy admissible paths and actual path velocities. It does not assert a triangle inequality for $\sqrt{\D_\Gamma^g}$, convergence of finite neural optimization, or universal approximation for a chosen constrained architecture. The transport bound concerns the population critic optimum and does not certify its empirical training loss.

The experimental scope also matters. A controlled flow baseline at a stated budget does not reproduce every author-budget RFM result. Sample discrepancies cannot be compared numerically with published NLL tables. The objective controls reveal substantial optimization failures: good initialized mixtures can degrade, and the matched direct-MMD MLP is stronger at the studied setting. Its evaluation-aligned kernel is itself an advantage, so this does not establish universal MMD superiority. The mesh surrogate, radial maps, finite-difference velocities, and global path-selection gradients need further analysis. Real large EEG covariance experiments are not completed, and modern one-step competitors require direct evaluation before any competitive advantage is claimed.

\section{Conclusion}
Riemannian midpoint learning supplies a geometric distribution discrepancy and a direct generator training rule. The population criterion identifies endpoint laws and controls Wasserstein-1 distance, with simulation-free closed-form interpolation where available. The controlled experiments do not yet establish a practical benefit of midpoint optimization over strong direct objectives: initialization, critic approximation, and gradient estimation remain consequential. The validated population theory and these empirical limitations motivate improved optimization and faithful comparisons with established one-step manifold methods.
